% ==============================================================================
% Resumo na língua vernácula (obrigatório, NBR 14724:2024 e NBR 6028:2021)
% ==============================================================================
\begin{resumo}
\begin{instrucao}
\orientacao{Resumo}
O resumo apresenta, de forma concisa, o objetivo, o método, os resultados e as
conclusões do trabalho. Escreva-o por último, quando o texto estiver pronto.
Siga a NBR 6028:2021:
\begin{itemize}
  \item parágrafo único, com 150 a 500 palavras;
  \item frases afirmativas, com o verbo na terceira pessoa, como recomenda o
        item 4.1.4 da norma (``Este trabalho propõe\ldots'', ``Os resultados
        mostram\ldots''). A voz ativa deixou de ser exigida, mas costuma
        deixar o resumo mais direto;
  \item sem citações, fórmulas, siglas não explicadas, figuras ou tabelas;
  \item palavras-chave logo abaixo, separadas por ponto e vírgula e finalizadas
        por ponto. Prefira termos que um leitor usaria para encontrar o
        trabalho num catálogo (de três a cinco é a prática mais comum).
\end{itemize}
Um roteiro que funciona bem dedica uma ou duas frases ao contexto e ao
problema, outras ao objetivo e ao método e as últimas aos principais
resultados (com números) e à conclusão.
\end{instrucao}

O problema da mochila multidimensional consiste em selecionar um subconjunto de
itens de maior valor total sem violar a capacidade de vários recursos ao mesmo
tempo. Ele modela decisões como a alocação de orçamento entre projetos e a
seleção de tarefas em servidores com limites de processamento e memória. Por
ser NP-difícil, instâncias grandes não costumam ser resolvidas na otimalidade
em tempo razoável por solucionadores de programação linear inteira. Este
trabalho propõe uma heurística de busca local iterada que combina uma solução
inicial gulosa, baseada em uma razão de eficiência ponderada pelas restrições,
com uma busca local de trocas e uma perturbação adaptativa. A heurística foi
implementada em C++ e comparada ao solucionador HiGHS e a um algoritmo
genético da literatura em 270 instâncias da OR-Library, com tempo limite de 60
segundos por instância. A heurística obteve desvio médio de 0,41\% em relação
às melhores soluções conhecidas, contra 1,87\% do solucionador no mesmo tempo, e
foi superior ao algoritmo genético em 82\% das instâncias. Os resultados
indicam que a perturbação adaptativa é o componente que mais contribui para o
desempenho, sobretudo nas instâncias com 500 itens.

\palavraschave{mochila multidimensional; busca local iterada; heurística;
otimização combinatória}
\end{resumo}
